\clearpage
\SourceLeaf{293}
\GaussVSourceUnnumbered
\pdfbookmark[1]{Göttingische gelehrte Anzeigen. 1829 October 12.}{anzeigen-1829-october-12}
\begin{center}
\rule{.70\linewidth}{.4pt}\\[.7\baselineskip]
Göttingische gelehrte Anzeigen. 1829 October 12.\\[.7\baselineskip]
\rule{.70\linewidth}{.4pt}
\end{center}

Am 28. September übergab Hofr. \textsc{Gauss} der Königl. Societät eine Vorlesung:
\begin{center}
\textit{Principia generalia theoriae figurae fluidorum in statu aequilibrii,}
\end{center}
\noindent
von welcher wir hier Bericht abzustatten haben. Ihr Gegenstand gehört in dasjenige
Gebiet der mathematischen Physik, welches \textsc{Laplace} durch seine beiden in den
Jahren 1806 und 1807 erschienenen, in dieser Wissenschaft Epoche machenden
Abhandlungen, \textit{Théorie de l'action capillaire}, und \textit{Supplément à la théorie
de l'action capillaire} auf eine so glänzende Art eröffnet hat. Zur Erklärung der
Gesetze, nach welchen wir die Himmelskörper sich bewegen sehen, ist die Annahme
einer allgemeinen gegenseitigen Anziehung alles Materiellen, deren Stärke dem
Quadrate der Entfernung umgekehrt proportional ist, sowohl nothwendig, als
zureichend. Die Schwere der Körper auf der Oberfläche der Erde ist gleichfalls
nichts weiter, als eine Wirkung dieser allgemeinen Kraft. Allein die mannigfaltigen
Erscheinungen, welche Körper in der Berührung darbieten, das Aufsteigen einiger
Flüssigkeiten in sehr engen Röhren, das Sinken anderer, die Adhäsion der Flüssigkeiten
an einigen festen Körpern, ihre Tropfengestalt oberhalb anderer, die Cohäsion
u. s. w. lassen sich aus einer jenes Gesetz befolgenden Anziehung nicht erklären,
da eine richtig geführte Rechnung leicht zeigt, dass ein einzelner Körper, dessen
Dimensionen gegen die der ganzen Erde verschwinden, gegen einen wo immer
befindlichen Punkt vermöge jenes Gesetzes
\clearpage
\SourceLeaf{294}
\GaussVSourceFolio{288}{L}{ANZEIGEN.}
\noindent
nur eine gegen die Schwere unmerkliche Anziehung ausüben kann. Man ist daher
genöthigt, das Gesetz der allgemeinen Anziehung für sehr kleine Entfernungen zu
modificiren, oder, was dasselbe ist, neben jener dem Quadrate der Entfernung
umgekehrt proportionalen Anziehung noch eine andere anzunehmen, deren eigentliches
Gesetz genau auszudrücken uns zwar die Mittel fehlen, die aber, wie die Erscheinungen
lehren, in jeder für uns messbaren Entfernung unmessbar klein, in unmessbar
kleinen Entfernungen hingegen nicht bloss merklich, sondern sogar überaus gross
sein muss. Man nennt diese Kraft Molecularanziehung, eine Benennung, die freilich
eigentlich nichts Bezeichnendes hat, wogegen der von einigen deutschen Physikern
gebrauchte Ausdruck Flächenkraft etwas anderes bezeichnet, als was hier bezeichnet
werden soll.

\textsc{Laplace} hat zuerst diese Vorstellung von der bei jenen Phänomenen thätigen
Kraft in Beziehung auf die Gleichgewichtsfigur der liquiden Flüssigkeiten dem
Calcül unterworfen, durch eine schöne Analyse die in jedem Punkte der Oberfläche
der Flüssigkeit zum Gleichgewicht nothwendige Gleichung aufgefunden, und nicht
bloss die eigentlich sogenannten Capillar-Phänomene, sondern eine Menge anderer
damit verwandter Erscheinungen daraus erklärt. Diese Untersuchungen, durchgehends
durch eine überraschend genaue Übereinstimmung mit der Erfahrung bestätigt,
gehören zu den schönsten Arbeiten jenes grossen Geometers. Zwar hat es nicht an
Gegnern dieser Theorie gefehlt: man findet jedoch nicht, dass bisher etwas irgend
erhebliches dagegen vorgebracht wäre, und die Schwäche der bekannt gewordenen
Einwürfe, welche höchstens einige Nachlässigkeiten in der Darstellung, aber nicht
die Sache selbst treffen, ist leicht nachzuweisen, wie z. B. \textsc{Petit} die von
\textsc{Brunacci} gemachten Einwendungen siegreich wiederlegt hat.

Um so mehr muss man sich wundern, dass eine wirkliche und wesentliche
Mangelhaftigkeit an dieser Theorie bisher ganz übersehen ist. Das Wesen dieser
Theorie beruht nämlich, genau betrachtet, auf zwei Hauptsätzen. Der eine besteht
in der vorhin erwähnten Gleichung, welche in jedem Punkte der freien Oberfläche
der Flüssigkeit beim Gleichgewicht Statt haben muss, und deren Begründung, wie
\textsc{Laplace} sie gegeben hat, nichts wesentliches zu wünschen übrig lässt.
Diese Gleichung ist eine partielle Differentialgleichung, die für sich allein die
Gestalt der Oberfläche nicht vollständig bestimmen kann, da ihre Integration,
wenn sie allgemein möglich wäre, noch zwei arbiträre Functionen einführen würde,
\clearpage
\SourceLeaf{295}
\GaussVSourceFolio{289}{R}{PRINCIPIA GENERALIA THEORIAE FIGURAE FLUIDORUM IN STATU AEQUILIBRII.}
\GaussVSourceGathering{37}
\noindent
deren Bestimmung anderswoher entlehnt werden muss. Die Stelle dieses Erfordernisses
vertritt nun der zweite Hauptsatz, nach welchem im Zustande des Gleichgewichts
die freie Oberfläche der Flüssigkeit da, wo sie durch das Gefäss begrenzt wird,
mit der Wand des Gefässes einen bestimmten constanten Winkel machen muss, der
bloss von dem Verhältniss der Molecularanziehungen abhängt, welche die Theile der
Flüssigkeit einerseits von einander und andererseits von den Theilen des Gefässes
erleiden. Dieser höchst wichtige Satz, ohne welchen die Theorie nur zur Hälfte
vollendet sein würde, gehört gleichfalls \textsc{Laplace} an, und ist in dessen
Theorie von Anfang bis zu Ende verwebt: allein umsonst sucht man in beiden
angeführten Schriften eine befriedigende bloss auf die Natur der Molecularanziehung
gestützte Begründung desselben. Was in der ersten Abhandlung S. 5 oben vorkommt,
setzt, was bewiesen werden sollte, schon voraus, und die Rechnungen in derselben
Abhandlung S. 44 u. f. führen zu gar keinem Resultate. Was sonst noch hierüber zu
sagen ist, muss hier der Kürze wegen übergangen, und demnächst in der vorliegenden
Abhandlung selbst nachgesehen werden.

Dieser Umstand ist eine von den Veranlassungen gewesen, die den Hofr.
\textsc{Gauss} bewogen haben, diese Untersuchung von neuem aufzunehmen, und zwar
auf einem eigenthümlichen Wege, der von dem von \textsc{Laplace} benutzten gänzlich
verschieden ist, wenn gleich jener und diese von einerlei Grundvoraussetzung in
Beziehung auf die Natur der Molecularanziehung ausgehen, und am Ende zu einerlei
Ziele führen. \textsc{Laplace} hat die erwähnte Gleichung für das Gleichgewicht
auf eine doppelte Art begründet; in der ersten Abhandlung mit Hülfe des Princips
des Gleichgewichts in unendlich engen Canälen; in der zweiten vermittelst des
Satzes, dass die Gesammtkraft, welche auf irgend einen Punkt der freien Oberfläche
der Flüssigkeit wirkt, beim Gleichgewicht auf die Oberfläche senkrecht ist. Die
gegenwärtige neue Bearbeitung der Theorie der Gleichgewichtsgestalt der Flüssigkeiten
geht dagegen von dem Princip der virtuellen Bewegungen aus.

Wir würden die dieser Anzeige gesetzten Grenzen weit überschreiten müssen,
wenn wir hier dem Gange der Untersuchungen im Einzelnen folgen wollten. Aber
verweilen müssen wir bei einem neuen Theorem, welches einen Hauptabschnitt in
denselben macht, und in einer einzigen Formel die Auflösung der Aufgabe in
grösster Einfachheit und Klarheit darstellt. Es ist folgendes:

Wenn man durch $s$ das Volumen der Flüssigkeit, durch $h$ die Höhe ihres
Schwerpunkts über einer beliebigen horizontalen Ebene, durch $T$ den Inhalt
\clearpage
\SourceLeaf{296}
\GaussVSourceFolio{290}{L}{ANZEIGEN.}
\noindent
desjenigen Theils der Oberfläche der Flüssigkeit, welche das Gefäss berührt,
und durch $U$ den Inhalt des andern (freien) Theils dieser Oberfläche bezeichnet:
so ist im Zustande des Gleichgewichts das Aggregat
\[
sh+(\alpha\alpha-2\beta\beta)T+\alpha\alpha U
\]
ein Minimum, wo $\alpha\alpha$ und $\beta\beta$ gewisse Constanten bedeuten,
welche von dem Verhältniss der Schwere zu der Intensität der Molecularanziehung
der Theile der Flüssigkeit gegen einander und der Theile des Gefässes gegen
die Flüssigkeit abhangen.

Wir sehen hier also, als die Frucht einer schwierigen und subtilen Untersuchung
einen Ausdruck für das Gesetz des Gleichgewichts hervorgehen, der, selbst dem
gemeinen Verstande begreiflich, die Vermittlung des Conflicts zwischen den
verschiedenen hier ins Spiel tretenden Kräften klar vor Augen legt. Wäre die
Schwere die einzige wirkende Kraft, so würde beim Gleichgewicht der Schwerpunkt
der ganzen Flüssigkeit so tief wie möglich liegen, also $h$ ein Minimum sein
müssen. Setzt man hingegen die Schwere und die Anziehung des Gefässes ganz bei
Seite, so dass bloss die gegenseitige Anziehung der Theile der Flüssigkeit selbst
in Betracht kommt, so muss diese eine sphärische Gestalt annehmen, also $T+U$
ein Minimum sein. Wäre endlich weder Schwere noch gegenseitige Anziehung der
Flüssigkeitstheile vorhanden, so würde die Flüssigkeit sich über die ganze
Oberfläche des Gefässes verbreiten, also $T$ ein Maximum, oder $-T$ ein Minimum
sein müssen. Man findet es begreiflich, dass beim Zusammenwirken der drei Kräfte
ein aus jenen drei Grössen Zusammengesetztes ein Kleinstes werden soll, wiewohl
sich von selbst versteht, dass die eigentliche feste Begründung jenes Lehrsatzes
nur auf die vollständigen strengen mathematischen Schlussreihen gestützt werden
kann, die von der Natur der Molecularanziehung wesentlich abhängig sind.

Mit grosser Leichtigkeit leitet man aus dieser Formel die Erscheinungen des
Steigens oder Fallens der Flüssigkeit in Haarröhrchen mit verticalen innern
Seitenwänden ab: hier beschränken wir uns auf eine kurze Andeutung und auf den
Fall, wo die Weite des Haarröhrchens gegen die Weite des Gefässes an der Oberfläche
der Flüssigkeit, in welche jenes eingetaucht ist, als verschwindend betrachtet
werden kann, mithin das Fallen oder Steigen der Flüssigkeit im Gefässe, welches
mit dem Steigen oder Fallen im Haarröhrchen zusammenhängt,
\clearpage
\SourceLeaf{297}
\GaussVSourceFolio{291}{R}{PRINCIPIA GENERALIA THEORIAE FIGURAE FLUIDORUM IN STATU AEQUILIBRII.}
\GaussVSourceGathering{37*}
\noindent
vernachlässigt werden darf. Da dies obige Aggregat seinen kleinsten Werth für
die Gleichgewichtsgestalt der Flüssigkeit hat, unter \textit{allen} anderen Gestalten,
welche diese annehmen kann, so ist jener Werth auch der kleinste unter denen,
die bei bloss veränderter Höhe der Flüssigkeit im Haarröhrchen ohne Veränderung
der Figur der freien Oberfläche in derselben hervorgehen. Bezeichnet man nun
mit $h^0,T^0,U^0$ die Werthe von $h,T,U$, wenn die Flüssigkeit im Haarröhrchen
eben so hoch steht, wie im Gefässe (oder genau, wenn sie im Mittel eben so
hoch steht, insofern die freie Oberfläche in jenem keine horizontale Ebene
bildet); ferner mit $a$ und $b$ den Flächeninhalt und den Umfang eines innern
Querschnitts des Haarröhrchens, so wird man allgemein, der (mittlern) Höhe im
Haarröhrchen über dem Niveau im Gefässe $=z$ entsprechend, haben
\[
\begin{aligned}
sh&=sh^0+\tfrac12azz\\
T&=T^0+bz\\
U&=U^0
\end{aligned}
\]
also das obige Aggregat
\[
=sh^0+(\alpha\alpha-2\beta\beta)T^0+\alpha\alpha U^0
 +\tfrac12azz+(\alpha\alpha-2\beta\beta)bz
\]
welches offenbar ein Minimum wird für
\[
z=\frac{(2\beta\beta-\alpha\alpha)b}{a}
\]
Die Flüssigkeit wird also im Haarröhrchen höher oder tiefer stehen, als im
Gefässe, je nachdem $\beta\beta$ grösser oder kleiner ist als $\frac12\alpha\alpha$,
und der Unterschied des Standes, bei bestimmter Flüssigkeit und Haarröhrchen
von bestimmter Materie wird der Peripherie des innern Querschnitts direct und
dem Flächeninhalt desselben verkehrt proportional sein.

Von höherer Wichtigkeit, als die Erledigung dieser zwar vorzüglich in die
Augen fallenden, aber doch nur ganz speciellen Erscheinung, ist die
\textit{allgemeine} Entwickelung der Folgen des obigen Lehrsatzes. Es ist klar,
dass dieses Geschäft der Variationsrechnung angehört, aber einem Theil derselben,
der bisher noch wenig oder gar nicht bearbeitet ist, wo nemlich von der
Variation doppelter Integrale mit veränderlichen Grenzen die Rede ist.
In dieser Beziehung werden die hier geöffneten Wege auch ein rein mathematisches
Interesse darbieten: in dieser Anzeige können wir nur die Resultate bemerken.
Diese bestehen, erstlich,
\clearpage
\SourceLeaf{298}
\GaussVSourceFolio{292}{L}{ANZEIGEN.}
\noindent
in einer Gleichung für jeden Punkt der freien Oberfläche der Flüssigkeit,
welche gänzlich mit der oben erwähnten \textsc{Laplace}schen Gleichung übereinstimmt;
zweitens, für die Grenzen dieser Fläche, in der Gleichung
$\sin\frac12 i=\frac{\beta}{\alpha}$, wenn man durch $i$ den Neigungswinkel
zwischen den die freie Oberfläche der Flüssigkeit und die Oberfläche des
Gefässes berührenden Ebenen, und zwar ausserhalb der Flüssigkeit gemessen,
bezeichnet, welches gerade der zweite von \textsc{Laplace} ohne Begründung
gebrauchte, Hauptsatz ist.

In einer Anzeige, wie sie diese Blätter verstatten, konnten nur Hauptmomente
der vorliegenden Arbeit berührt werden; vieles andere, was diese umfasst,
übergehen wir hier ganz oder deuten es nur kurz an. Dahin gehören die
Modificationen, welche nothwendig werden, wenn ein Theil der Flüssigkeit als
ein Häutchen von unmessbar kleiner Dicke an der Gefässwand anliegt (diese
benetzt), welcher Fall eintritt, wenn $\beta$ grösser ist als $\alpha$, wo
die Gleichung $\sin\frac12 i=\frac{\beta}{\alpha}$ eine scheinbare Ungereimtheit
enthält; die Folgen des Umstandes, dass die Kenntniss des Gesetzes der
Molecularanziehung in unmessbar kleinen Distanzen (deren Begriff ganz verschieden
ist von verschwindenden oder von unendlich kleinen Distanzen) uns unzugänglich
bleibt; die Unterscheidung zwischen dem Zustand des wahren Gleichgewichts und
des wegen der Reibung an den Gefässwänden auch ausser demselben möglichen
Zustandes der Ruhe; die Betrachtung der Discontinuität der Gestalt der
Gefässwände (einwärts oder auswärts gehender Winkel) u. s. f.

Der eigentliche Zweck der vorliegenden Arbeit ging dahin, eine feste allgemeine
mathematische Begründung der Hauptprincipien dieser Lehre zu geben: es lag,
für jetzt, ausserhalb des Planes, specielle Phänomene zu erklären, worin
ohnehin \textsc{Laplace} schon so viel geleistet hat. Der hier gewählte Weg
ist nicht der einzige: der Verf. hat während der Ausarbeitung noch einen ganz
andern, nicht weniger merkwürdigen Weg zu demselben Ziele, und namentlich zu
dem schwer zugänglichen zweiten Hauptlehrsatz gefunden, dessen Darstellung,
so wie die weitere Ausdehnung der Principien auf die zusammengesetzten Fälle,
wo mehrere Flüssigkeiten in einem Gefässe, eingetauchte ganz oder zum Theil
freie feste Körper u. dergl. mit in Betracht kommen, er sich auf eine andere
Gelegenheit vorbehält.
